
\documentclass{article}
\usepackage{colortbl}
\usepackage{makecell}
\usepackage{multirow}
\usepackage{supertabular}

\begin{document}

\newcounter{utterance}

\centering \large Interaction Transcript for game `matchit\_ascii', experiment `similar\_grid\_2', episode 3 with claude{-}opus{-}4{-}8{-}azure.
\vspace{24pt}

{ \footnotesize  \setcounter{utterance}{1}
\setlength{\tabcolsep}{0pt}
\begin{supertabular}{c@{$\;$}|p{.15\linewidth}@{}p{.15\linewidth}p{.15\linewidth}p{.15\linewidth}p{.15\linewidth}p{.15\linewidth}}
    \# & $\;$A & \multicolumn{4}{c}{Game Master} & $\;\:$B\\
    \hline

    \theutterance \stepcounter{utterance}  
    & & \multicolumn{4}{p{0.6\linewidth}}{
        \cellcolor[rgb]{0.9,0.9,0.9}{
            \makecell[{{p{\linewidth}}}]{
                \texttt{\tiny{[P1$\langle$GM]}}
                \texttt{İşbirlikçi bir tahmin oyununa katılıyorsun.} \\
\texttt{Amaç, bu ızgara ile yalnızca benim görebildiğim diğer görselin aynı olup olmadığını anlamak.} \\
\texttt{Önce sen kendi ızgaranı tarif et. Ardından ben de kendi görselimin tarifini vereceğim ve ızgaraların aynı olup olmadığını anlamak için birbirimize sorular sorabileceğiz.  Şimdi ızgaranın kısa bir tarifini "\textbackslash{}"AÇIKLAMA:\textbackslash{}"" ile başlayacak şekilde yaz ve ardından açıklamanı ekle. Başka hiçbir şey ekleme. İşte senin ızgaran:} \\
\texttt{X X X X X} \\
\texttt{▢ ▢ ▢ ▢ X} \\
\texttt{▢ ▢ X X X} \\
\texttt{▢ ▢ ▢ ▢ X} \\
\texttt{▢ ▢ ▢ ▢ X} \\
            }
        }
    }
    & \\ \\

    \theutterance \stepcounter{utterance}  
    & \multicolumn{4}{p{0.6\linewidth}}{
        \cellcolor[rgb]{0.9,0.9,0.9}{
            \makecell[{{p{\linewidth}}}]{
                \texttt{\tiny{[P1$\rangle$GM]}}
                \texttt{AÇIKLAMA: 5x5'lik bir ızgara. İlk satır tamamen dolu (X X X X X). İkinci satırda sadece son hücre dolu. Üçüncü satırda son üç hücre dolu (X X X). Dördüncü ve beşinci satırlarda sadece son hücre dolu. Yani sağ kenar boyunca dikey bir çizgi, üstte tam bir yatay çizgi ve ortada sağa doğru uzanan kısa bir yatay çizgi var.} \\
            }
        }
    }
    & & \\ \\

    \theutterance \stepcounter{utterance}  
    & & & \multicolumn{2}{p{0.3\linewidth}}{
        \cellcolor[rgb]{0.9,0.9,0.9}{
            \makecell[{{p{\linewidth}}}]{
                \texttt{\tiny{[GM$|$GM]}}
                \texttt{continue} \\
            }
        }
    }
    & & \\ \\

    \theutterance \stepcounter{utterance}  
    & & \multicolumn{4}{p{0.6\linewidth}}{
        \cellcolor[rgb]{0.9,0.9,0.9}{
            \makecell[{{p{\linewidth}}}]{
                \texttt{\tiny{[GM$\rangle$P2]}}
                \texttt{İşbirlikçi bir tahmin oyununa katılıyorsun.} \\
\texttt{Amaç, bu ızgara ile yalnızca benim görebildiğim diğer görselin aynı olup olmadığını anlamak.} \\
\texttt{Önce sen kendi ızgaranı tarif et. Ardından ben de kendi görselimin tarifini vereceğim ve ızgaraların aynı olup olmadığını anlamak için birbirimize sorular sorabileceğiz.  Şimdi ızgaranın kısa bir tarifini "\textbackslash{}"AÇIKLAMA:\textbackslash{}"" ile başlayacak şekilde yaz ve ardından açıklamanı ekle. Başka hiçbir şey ekleme. İşte senin ızgaran:} \\
\texttt{X X X X X} \\
\texttt{▢ ▢ ▢ ▢ X} \\
\texttt{▢ X X X X} \\
\texttt{▢ ▢ ▢ ▢ X} \\
\texttt{▢ ▢ X ▢ X} \\
            }
        }
    }
    & \\ \\

    \theutterance \stepcounter{utterance}  
    & & & \multicolumn{4}{p{0.6\linewidth}}{
        \cellcolor[rgb]{0.9,0.9,0.9}{
            \makecell[{{p{\linewidth}}}]{
                \texttt{\tiny{[GM$\langle$P2]}}
                \texttt{AÇIKLAMA: 5x5'lik bir ızgaram var. Satır satır anlatayım:} \\
\texttt{{-} 1. satır: Tüm hücreler dolu (X X X X X)} \\
\texttt{{-} 2. satır: İlk dört hücre boş, son hücre dolu (▢ ▢ ▢ ▢ X)} \\
\texttt{{-} 3. satır: İlk hücre boş, kalan dördü dolu (▢ X X X X)} \\
\texttt{{-} 4. satır: İlk dört hücre boş, son hücre dolu (▢ ▢ ▢ ▢ X)} \\
\texttt{{-} 5. satır: İlk iki hücre boş, üçüncü dolu, dördüncü boş, son hücre dolu (▢ ▢ X ▢ X)} \\
\\ 
\texttt{Toplamda sağ kenar (son sütun) tamamen dolu ve üstte bir çizgi oluşturuyor; bir tür yılankavi/zikzak yol gibi görünüyor.} \\
            }
        }
    }
     \\ \\

    \theutterance \stepcounter{utterance}  
    & & & \multicolumn{2}{p{0.3\linewidth}}{
        \cellcolor[rgb]{0.9,0.9,0.9}{
            \makecell[{{p{\linewidth}}}]{
                \texttt{\tiny{[GM$|$GM]}}
                \texttt{continue} \\
            }
        }
    }
    & & \\ \\

    \theutterance \stepcounter{utterance}  
    & & \multicolumn{4}{p{0.6\linewidth}}{
        \cellcolor[rgb]{0.9,0.9,0.9}{
            \makecell[{{p{\linewidth}}}]{
                \texttt{\tiny{[GM$\rangle$P2]}}
                \texttt{Bu benimAÇIKLAMA: 5x5'lik bir ızgara. İlk satır tamamen dolu (X X X X X). İkinci satırda sadece son hücre dolu. Üçüncü satırda son üç hücre dolu (X X X). Dördüncü ve beşinci satırlarda sadece son hücre dolu. Yani sağ kenar boyunca dikey bir çizgi, üstte tam bir yatay çizgi ve ortada sağa doğru uzanan kısa bir yatay çizgi var.} \\
\texttt{Şimdi, ızgaramın ızgarandan farklı olabilecek yeni özelliklerini öğrenmek için bir soru sor. Soruna "SORU:" ile başla ve başka hiçbir şey ekleme.} \\
            }
        }
    }
    & \\ \\

    \theutterance \stepcounter{utterance}  
    & & & \multicolumn{4}{p{0.6\linewidth}}{
        \cellcolor[rgb]{0.9,0.9,0.9}{
            \makecell[{{p{\linewidth}}}]{
                \texttt{\tiny{[GM$\langle$P2]}}
                \texttt{SORU: Üçüncü satırında ikinci hücre (soldan) dolu mu, boş mu?} \\
            }
        }
    }
     \\ \\

    \theutterance \stepcounter{utterance}  
    & & & \multicolumn{2}{p{0.3\linewidth}}{
        \cellcolor[rgb]{0.9,0.9,0.9}{
            \makecell[{{p{\linewidth}}}]{
                \texttt{\tiny{[GM$|$GM]}}
                \texttt{continue} \\
            }
        }
    }
    & & \\ \\

    \theutterance \stepcounter{utterance}  
    & & \multicolumn{4}{p{0.6\linewidth}}{
        \cellcolor[rgb]{0.9,0.9,0.9}{
            \makecell[{{p{\linewidth}}}]{
                \texttt{\tiny{[P1$\langle$GM]}}
                \texttt{Bu benimAÇIKLAMA: 5x5'lik bir ızgaram var. Satır satır anlatayım:} \\
\texttt{{-} 1. satır: Tüm hücreler dolu (X X X X X)} \\
\texttt{{-} 2. satır: İlk dört hücre boş, son hücre dolu (▢ ▢ ▢ ▢ X)} \\
\texttt{{-} 3. satır: İlk hücre boş, kalan dördü dolu (▢ X X X X)} \\
\texttt{{-} 4. satır: İlk dört hücre boş, son hücre dolu (▢ ▢ ▢ ▢ X)} \\
\texttt{{-} 5. satır: İlk iki hücre boş, üçüncü dolu, dördüncü boş, son hücre dolu (▢ ▢ X ▢ X)} \\
\\ 
\texttt{Toplamda sağ kenar (son sütun) tamamen dolu ve üstte bir çizgi oluşturuyor; bir tür yılankavi/zikzak yol gibi görünüyor.} \\
\texttt{SORU: Üçüncü satırında ikinci hücre (soldan) dolu mu, boş mu?Yanıtına "CEVAP:" ile başla ve başka hiçbir şey ekleme.} \\
            }
        }
    }
    & \\ \\

    \theutterance \stepcounter{utterance}  
    & \multicolumn{4}{p{0.6\linewidth}}{
        \cellcolor[rgb]{0.9,0.9,0.9}{
            \makecell[{{p{\linewidth}}}]{
                \texttt{\tiny{[P1$\rangle$GM]}}
                \texttt{CEVAP: Benim üçüncü satırımda ikinci hücre (soldan) boş.} \\
            }
        }
    }
    & & \\ \\

    \theutterance \stepcounter{utterance}  
    & & & \multicolumn{2}{p{0.3\linewidth}}{
        \cellcolor[rgb]{0.9,0.9,0.9}{
            \makecell[{{p{\linewidth}}}]{
                \texttt{\tiny{[GM$|$GM]}}
                \texttt{continue} \\
            }
        }
    }
    & & \\ \\

    \theutterance \stepcounter{utterance}  
    & & \multicolumn{4}{p{0.6\linewidth}}{
        \cellcolor[rgb]{0.9,0.9,0.9}{
            \makecell[{{p{\linewidth}}}]{
                \texttt{\tiny{[P1$\langle$GM]}}
                \texttt{Şimdi, ızgaramın ızgarandan farklı olabilecek yeni özelliklerini öğrenmek için bir soru sor. Soruna "SORU:" ile başla ve başka hiçbir şey ekleme.} \\
            }
        }
    }
    & \\ \\

    \theutterance \stepcounter{utterance}  
    & \multicolumn{4}{p{0.6\linewidth}}{
        \cellcolor[rgb]{0.9,0.9,0.9}{
            \makecell[{{p{\linewidth}}}]{
                \texttt{\tiny{[P1$\rangle$GM]}}
                \texttt{SORU: Beşinci satırında üçüncü hücre (soldan) dolu mu, boş mu?} \\
            }
        }
    }
    & & \\ \\

    \theutterance \stepcounter{utterance}  
    & & & \multicolumn{2}{p{0.3\linewidth}}{
        \cellcolor[rgb]{0.9,0.9,0.9}{
            \makecell[{{p{\linewidth}}}]{
                \texttt{\tiny{[GM$|$GM]}}
                \texttt{continue} \\
            }
        }
    }
    & & \\ \\

    \theutterance \stepcounter{utterance}  
    & & \multicolumn{4}{p{0.6\linewidth}}{
        \cellcolor[rgb]{0.9,0.9,0.9}{
            \makecell[{{p{\linewidth}}}]{
                \texttt{\tiny{[GM$\rangle$P2]}}
                \texttt{CEVAP: Benim üçüncü satırımda ikinci hücre (soldan) boş.} \\
\texttt{SORU: Beşinci satırında üçüncü hücre (soldan) dolu mu, boş mu?Yanıtına "CEVAP:" ile başla ve başka hiçbir şey ekleme.} \\
            }
        }
    }
    & \\ \\

    \theutterance \stepcounter{utterance}  
    & & & \multicolumn{4}{p{0.6\linewidth}}{
        \cellcolor[rgb]{0.9,0.9,0.9}{
            \makecell[{{p{\linewidth}}}]{
                \texttt{\tiny{[GM$\langle$P2]}}
                \texttt{CEVAP: Benim beşinci satırımda üçüncü hücre (soldan) dolu. Ayrıca belirtmeliyim ki benim üçüncü satırımda ikinci hücre dolu, seninkinde ise boş. Bu nedenle ızgaralarımız farklı.} \\
            }
        }
    }
     \\ \\

    \theutterance \stepcounter{utterance}  
    & & & \multicolumn{2}{p{0.3\linewidth}}{
        \cellcolor[rgb]{0.9,0.9,0.9}{
            \makecell[{{p{\linewidth}}}]{
                \texttt{\tiny{[GM$|$GM]}}
                \texttt{continue} \\
            }
        }
    }
    & & \\ \\

    \theutterance \stepcounter{utterance}  
    & & \multicolumn{4}{p{0.6\linewidth}}{
        \cellcolor[rgb]{0.9,0.9,0.9}{
            \makecell[{{p{\linewidth}}}]{
                \texttt{\tiny{[GM$\rangle$P2]}}
                \texttt{Şimdi, ızgaramın ızgarandan farklı olabilecek yeni özelliklerini öğrenmek için bir soru sor. Soruna "SORU:" ile başla ve başka hiçbir şey ekleme.} \\
            }
        }
    }
    & \\ \\

    \theutterance \stepcounter{utterance}  
    & & & \multicolumn{4}{p{0.6\linewidth}}{
        \cellcolor[rgb]{0.9,0.9,0.9}{
            \makecell[{{p{\linewidth}}}]{
                \texttt{\tiny{[GM$\langle$P2]}}
                \texttt{SORU: Beşinci satırında üçüncü hücre (soldan) dolu mu, boş mu?} \\
            }
        }
    }
     \\ \\

    \theutterance \stepcounter{utterance}  
    & & & \multicolumn{2}{p{0.3\linewidth}}{
        \cellcolor[rgb]{0.9,0.9,0.9}{
            \makecell[{{p{\linewidth}}}]{
                \texttt{\tiny{[GM$|$GM]}}
                \texttt{continue} \\
            }
        }
    }
    & & \\ \\

    \theutterance \stepcounter{utterance}  
    & & \multicolumn{4}{p{0.6\linewidth}}{
        \cellcolor[rgb]{0.9,0.9,0.9}{
            \makecell[{{p{\linewidth}}}]{
                \texttt{\tiny{[P1$\langle$GM]}}
                \texttt{CEVAP: Benim beşinci satırımda üçüncü hücre (soldan) dolu. Ayrıca belirtmeliyim ki benim üçüncü satırımda ikinci hücre dolu, seninkinde ise boş. Bu nedenle ızgaralarımız farklı.} \\
\texttt{SORU: Beşinci satırında üçüncü hücre (soldan) dolu mu, boş mu?Yanıtına "CEVAP:" ile başla ve başka hiçbir şey ekleme.} \\
            }
        }
    }
    & \\ \\

    \theutterance \stepcounter{utterance}  
    & \multicolumn{4}{p{0.6\linewidth}}{
        \cellcolor[rgb]{0.9,0.9,0.9}{
            \makecell[{{p{\linewidth}}}]{
                \texttt{\tiny{[P1$\rangle$GM]}}
                \texttt{CEVAP: Benim beşinci satırımda üçüncü hücre (soldan) boş. Ayrıca belirttiğin gibi ızgaralarımız farklı.} \\
            }
        }
    }
    & & \\ \\

    \theutterance \stepcounter{utterance}  
    & & & \multicolumn{2}{p{0.3\linewidth}}{
        \cellcolor[rgb]{0.9,0.9,0.9}{
            \makecell[{{p{\linewidth}}}]{
                \texttt{\tiny{[GM$|$GM]}}
                \texttt{continue} \\
            }
        }
    }
    & & \\ \\

    \theutterance \stepcounter{utterance}  
    & & \multicolumn{4}{p{0.6\linewidth}}{
        \cellcolor[rgb]{0.9,0.9,0.9}{
            \makecell[{{p{\linewidth}}}]{
                \texttt{\tiny{[P1$\langle$GM]}}
                \texttt{Şimdi, ızgaramın ızgarandan farklı olabilecek yeni özelliklerini öğrenmek için bir soru sor. Soruna "SORU:" ile başla ve başka hiçbir şey ekleme.} \\
            }
        }
    }
    & \\ \\

    \theutterance \stepcounter{utterance}  
    & \multicolumn{4}{p{0.6\linewidth}}{
        \cellcolor[rgb]{0.9,0.9,0.9}{
            \makecell[{{p{\linewidth}}}]{
                \texttt{\tiny{[P1$\rangle$GM]}}
                \texttt{SORU: Üçüncü satırında ikinci hücre (soldan) hariç, üçüncü, dördüncü ve beşinci hücreler dolu mu?} \\
            }
        }
    }
    & & \\ \\

    \theutterance \stepcounter{utterance}  
    & & & \multicolumn{2}{p{0.3\linewidth}}{
        \cellcolor[rgb]{0.9,0.9,0.9}{
            \makecell[{{p{\linewidth}}}]{
                \texttt{\tiny{[GM$|$GM]}}
                \texttt{continue} \\
            }
        }
    }
    & & \\ \\

    \theutterance \stepcounter{utterance}  
    & & \multicolumn{4}{p{0.6\linewidth}}{
        \cellcolor[rgb]{0.9,0.9,0.9}{
            \makecell[{{p{\linewidth}}}]{
                \texttt{\tiny{[GM$\rangle$P2]}}
                \texttt{CEVAP: Benim beşinci satırımda üçüncü hücre (soldan) boş. Ayrıca belirttiğin gibi ızgaralarımız farklı.} \\
\texttt{SORU: Üçüncü satırında ikinci hücre (soldan) hariç, üçüncü, dördüncü ve beşinci hücreler dolu mu?Yanıtına "CEVAP:" ile başla ve başka hiçbir şey ekleme.} \\
            }
        }
    }
    & \\ \\

    \theutterance \stepcounter{utterance}  
    & & & \multicolumn{4}{p{0.6\linewidth}}{
        \cellcolor[rgb]{0.9,0.9,0.9}{
            \makecell[{{p{\linewidth}}}]{
                \texttt{\tiny{[GM$\langle$P2]}}
                \texttt{CEVAP: Evet, benim üçüncü satırımda üçüncü, dördüncü ve beşinci hücreler dolu. Zaten daha önce belirlediğimiz gibi ızgaralarımız farklı, çünkü senin beşinci satırında üçüncü hücre boşken benimkinde dolu ve senin üçüncü satırında ikinci hücre boşken benimkinde dolu.} \\
            }
        }
    }
     \\ \\

    \theutterance \stepcounter{utterance}  
    & & & \multicolumn{2}{p{0.3\linewidth}}{
        \cellcolor[rgb]{0.9,0.9,0.9}{
            \makecell[{{p{\linewidth}}}]{
                \texttt{\tiny{[GM$|$GM]}}
                \texttt{continue} \\
            }
        }
    }
    & & \\ \\

    \theutterance \stepcounter{utterance}  
    & & \multicolumn{4}{p{0.6\linewidth}}{
        \cellcolor[rgb]{0.9,0.9,0.9}{
            \makecell[{{p{\linewidth}}}]{
                \texttt{\tiny{[GM$\rangle$P2]}}
                \texttt{Şimdi, ızgaramın ızgarandan farklı olabilecek yeni özelliklerini öğrenmek için bir soru sor. Soruna "SORU:" ile başla ve başka hiçbir şey ekleme.} \\
            }
        }
    }
    & \\ \\

    \theutterance \stepcounter{utterance}  
    & & & \multicolumn{4}{p{0.6\linewidth}}{
        \cellcolor[rgb]{0.9,0.9,0.9}{
            \makecell[{{p{\linewidth}}}]{
                \texttt{\tiny{[GM$\langle$P2]}}
                \texttt{SORU: İkinci satırında sadece son hücre mi dolu, yoksa başka dolu hücre var mı?} \\
            }
        }
    }
     \\ \\

    \theutterance \stepcounter{utterance}  
    & & & \multicolumn{2}{p{0.3\linewidth}}{
        \cellcolor[rgb]{0.9,0.9,0.9}{
            \makecell[{{p{\linewidth}}}]{
                \texttt{\tiny{[GM$|$GM]}}
                \texttt{continue} \\
            }
        }
    }
    & & \\ \\

    \theutterance \stepcounter{utterance}  
    & & \multicolumn{4}{p{0.6\linewidth}}{
        \cellcolor[rgb]{0.9,0.9,0.9}{
            \makecell[{{p{\linewidth}}}]{
                \texttt{\tiny{[P1$\langle$GM]}}
                \texttt{CEVAP: Evet, benim üçüncü satırımda üçüncü, dördüncü ve beşinci hücreler dolu. Zaten daha önce belirlediğimiz gibi ızgaralarımız farklı, çünkü senin beşinci satırında üçüncü hücre boşken benimkinde dolu ve senin üçüncü satırında ikinci hücre boşken benimkinde dolu.} \\
\texttt{SORU: İkinci satırında sadece son hücre mi dolu, yoksa başka dolu hücre var mı?Yanıtına "CEVAP:" ile başla ve başka hiçbir şey ekleme.} \\
            }
        }
    }
    & \\ \\

    \theutterance \stepcounter{utterance}  
    & \multicolumn{4}{p{0.6\linewidth}}{
        \cellcolor[rgb]{0.9,0.9,0.9}{
            \makecell[{{p{\linewidth}}}]{
                \texttt{\tiny{[P1$\rangle$GM]}}
                \texttt{CEVAP: Benim ikinci satırımda sadece son hücre dolu, başka dolu hücre yok.} \\
            }
        }
    }
    & & \\ \\

    \theutterance \stepcounter{utterance}  
    & & & \multicolumn{2}{p{0.3\linewidth}}{
        \cellcolor[rgb]{0.9,0.9,0.9}{
            \makecell[{{p{\linewidth}}}]{
                \texttt{\tiny{[GM$|$GM]}}
                \texttt{continue} \\
            }
        }
    }
    & & \\ \\

    \theutterance \stepcounter{utterance}  
    & & \multicolumn{4}{p{0.6\linewidth}}{
        \cellcolor[rgb]{0.9,0.9,0.9}{
            \makecell[{{p{\linewidth}}}]{
                \texttt{\tiny{[P1$\langle$GM]}}
                \texttt{Şimdi, ızgaramın ızgarandan farklı olabilecek yeni özelliklerini öğrenmek için bir soru sor. Soruna "SORU:" ile başla ve başka hiçbir şey ekleme.} \\
            }
        }
    }
    & \\ \\

    \theutterance \stepcounter{utterance}  
    & \multicolumn{4}{p{0.6\linewidth}}{
        \cellcolor[rgb]{0.9,0.9,0.9}{
            \makecell[{{p{\linewidth}}}]{
                \texttt{\tiny{[P1$\rangle$GM]}}
                \texttt{SORU: Dördüncü satırında sadece son hücre mi dolu, yoksa başka dolu hücre var mı?} \\
            }
        }
    }
    & & \\ \\

    \theutterance \stepcounter{utterance}  
    & & & \multicolumn{2}{p{0.3\linewidth}}{
        \cellcolor[rgb]{0.9,0.9,0.9}{
            \makecell[{{p{\linewidth}}}]{
                \texttt{\tiny{[GM$|$GM]}}
                \texttt{continue} \\
            }
        }
    }
    & & \\ \\

    \theutterance \stepcounter{utterance}  
    & & \multicolumn{4}{p{0.6\linewidth}}{
        \cellcolor[rgb]{0.9,0.9,0.9}{
            \makecell[{{p{\linewidth}}}]{
                \texttt{\tiny{[GM$\rangle$P2]}}
                \texttt{CEVAP: Benim ikinci satırımda sadece son hücre dolu, başka dolu hücre yok.} \\
\texttt{SORU: Dördüncü satırında sadece son hücre mi dolu, yoksa başka dolu hücre var mı?Yanıtına "CEVAP:" ile başla ve başka hiçbir şey ekleme.} \\
            }
        }
    }
    & \\ \\

    \theutterance \stepcounter{utterance}  
    & & & \multicolumn{4}{p{0.6\linewidth}}{
        \cellcolor[rgb]{0.9,0.9,0.9}{
            \makecell[{{p{\linewidth}}}]{
                \texttt{\tiny{[GM$\langle$P2]}}
                \texttt{CEVAP: Benim dördüncü satırımda sadece son hücre dolu, başka dolu hücre yok. Bu satır bizde aynı. Ancak daha önce belirlediğimiz farklılıklar nedeniyle ızgaralarımız yine de farklı.} \\
            }
        }
    }
     \\ \\

    \theutterance \stepcounter{utterance}  
    & & & \multicolumn{2}{p{0.3\linewidth}}{
        \cellcolor[rgb]{0.9,0.9,0.9}{
            \makecell[{{p{\linewidth}}}]{
                \texttt{\tiny{[GM$|$GM]}}
                \texttt{continue} \\
            }
        }
    }
    & & \\ \\

    \theutterance \stepcounter{utterance}  
    & & \multicolumn{4}{p{0.6\linewidth}}{
        \cellcolor[rgb]{0.9,0.9,0.9}{
            \makecell[{{p{\linewidth}}}]{
                \texttt{\tiny{[GM$\rangle$P2]}}
                \texttt{Şimdi bir karara var. Ne düşünüyorsun: senin ızgaran ve diğer ızgara aynı mı tarif ediliyor? Aynı ızgara olduklarını düşünüyorsan "KARAR: aynı ızgara" yaz ya da farklı ızgaralar olduklarını düşünüyorsan "KARAR: farklı ızgaralar" yaz. Başka hiçbir şey ekleme.} \\
            }
        }
    }
    & \\ \\

    \theutterance \stepcounter{utterance}  
    & & & \multicolumn{4}{p{0.6\linewidth}}{
        \cellcolor[rgb]{0.9,0.9,0.9}{
            \makecell[{{p{\linewidth}}}]{
                \texttt{\tiny{[GM$\langle$P2]}}
                \texttt{KARAR: farklı ızgaralar} \\
            }
        }
    }
     \\ \\

    \theutterance \stepcounter{utterance}  
    & & & \multicolumn{2}{p{0.3\linewidth}}{
        \cellcolor[rgb]{0.9,0.9,0.9}{
            \makecell[{{p{\linewidth}}}]{
                \texttt{\tiny{[GM$|$GM]}}
                \texttt{continue} \\
            }
        }
    }
    & & \\ \\

    \theutterance \stepcounter{utterance}  
    & & & \multicolumn{2}{p{0.3\linewidth}}{
        \cellcolor[rgb]{0.9,0.9,0.9}{
            \makecell[{{p{\linewidth}}}]{
                \texttt{\tiny{[GM$|$GM]}}
                \texttt{success} \\
            }
        }
    }
    & & \\ \\

    \theutterance \stepcounter{utterance}  
    & & \multicolumn{4}{p{0.6\linewidth}}{
        \cellcolor[rgb]{0.9,0.9,0.9}{
            \makecell[{{p{\linewidth}}}]{
                \texttt{\tiny{[P1$\langle$GM]}}
                \texttt{CEVAP: Benim dördüncü satırımda sadece son hücre dolu, başka dolu hücre yok. Bu satır bizde aynı. Ancak daha önce belirlediğimiz farklılıklar nedeniyle ızgaralarımız yine de farklı.} \\
\texttt{Şimdi bir karara var. Ne düşünüyorsun: senin ızgaran ve diğer ızgara aynı mı tarif ediliyor? Aynı ızgara olduklarını düşünüyorsan "KARAR: aynı ızgara" yaz ya da farklı ızgaralar olduklarını düşünüyorsan "KARAR: farklı ızgaralar" yaz. Başka hiçbir şey ekleme.} \\
            }
        }
    }
    & \\ \\

    \theutterance \stepcounter{utterance}  
    & \multicolumn{4}{p{0.6\linewidth}}{
        \cellcolor[rgb]{0.9,0.9,0.9}{
            \makecell[{{p{\linewidth}}}]{
                \texttt{\tiny{[P1$\rangle$GM]}}
                \texttt{KARAR: farklı ızgaralar} \\
            }
        }
    }
    & & \\ \\

    \theutterance \stepcounter{utterance}  
    & & & \multicolumn{2}{p{0.3\linewidth}}{
        \cellcolor[rgb]{0.9,0.9,0.9}{
            \makecell[{{p{\linewidth}}}]{
                \texttt{\tiny{[GM$|$GM]}}
                \texttt{continue} \\
            }
        }
    }
    & & \\ \\

    \theutterance \stepcounter{utterance}  
    & & & \multicolumn{2}{p{0.3\linewidth}}{
        \cellcolor[rgb]{0.9,0.9,0.9}{
            \makecell[{{p{\linewidth}}}]{
                \texttt{\tiny{[GM$|$GM]}}
                \texttt{success} \\
            }
        }
    }
    & & \\ \\

\end{supertabular}
}

\end{document}
